\par\SourceLeaf{142}
In such circumstances the importance of the office of the chairman of a deliberative body becomes evident: he must distinguish the individual questions sharply, put each in its proper place, and submit them for deliberation and a vote in clear wording that excludes every ambiguity.

In the present case this task fell to F., as the prorector then in office. On 22 July he circulated the rescript, \textsc{Kritter}'s opinion and P.'s P. M. among all the professors, with a request which I give here in full in F.'s own words:
\begin{quote}
You are requested to state in writing whether you agree to P.'s entire proposal, or to express your views on the individual questions, namely 1) to what extent the increase in pensions seems just and advisable? 2) Whether an extension of the duration of the pension after the parents' death until the youngest child's \(20^{\text{th}}\) year is desired? 3) Whether the increase in annual contributions to 10 Thl. C. M. proposed by the Royal Government is approved, without condition; or subject to the condition in No. 2.
\end{quote}
\noindent\hspace*{\parindent}One sees that the last article of P.'s P. M. (above, p. [136]) is not mentioned by a single word. I myself now have no doubt that F. may have regarded it (perhaps not sufficiently appreciating the great importance of ‘this circumstance’) as an appendage already included in his No. 1: but I have just as little doubt that this failure to mention it, especially in contrast with the precisely logical form in which F. presents his third question, contributed greatly to pushing that main question into the background for many, perhaps for most, of those voting. I therefore thought it necessary to touch here on these incidental circumstances, insignificant in themselves, because they make more intelligible the otherwise so striking fact that of 41 voters not a single one entered into a discussion of that main question (unless one wishes to count E.'s vote on p. [139] as such); indeed, most voters did not mention it at all, and in fact only two voters mentioned it at all in an entirely unambiguous way.

These written deliberations (from 22 July to 6 August) are very voluminous, and some individual votes very detailed; but they turn almost exclusively on F.'s questions 2 and 3, which met with much opposition, especially the last; also on some new proposals introduced in the course of the deliberations, notably that of abolishing the existing freedom to join the widows' fund only later on paying double the contributions in arrears, a proposal which some voters supported and others emphatically rejected. I shall cite only a few of these votes which are more closely connected with my subject.

R. declares himself averse to all changes in general, but intends to accept the decisions of the majority. Although he does not mention P.'s Art. 3 at all, he does mention the underlying question of principle, on which he very firmly opposes the assertion made by P. in the earlier vote (above, p. [133], l. [21]).

He cannot, he says, convince himself of the correctness of the principle that, when the present widows' pension is increased, care must be taken that no necessity should arise later to reduce it and give future widows less than the present ones receive. This principle would certainly be correct if (as with other widows' funds) the fund's existence rested solely on the contributions. But since the professors' widows' fund is a benefaction, the contributions count for almost nothing, and participation in the benefaction is a condition of every professor's appointment, the far more natural principle is
\begin{quote}
Every widow receives the highest possible pension that the fund permits during her lifetime.
\end{quote}
\par\SourceLeaf{143}
On this principle no one has reason to complain, whether the fund rises or falls. If one says that future widows may receive less if the fund falls, he (R.) answers that present widows may receive less than future ones if the fund rises, and that this latter case is more probable than the former.

Thus R., as one sees, implicitly counters P.'s assertion that a future widow would suffer an injury if she received less than an earlier one had received, by replying that one could then, with exactly the same right, assert that a present widow would suffer an injury if she received less than a future one.

\textsc{Sch.} repeats the principle just cited (‘Every widow receives the highest possible, etc.’), adding: ‘H. H. R., it seems to me, has made this proposition evident, to which I refer. \textit{It also underlies No. 3 of P.'s proposal}’. After \textsc{Sch.} has also emphasised that the fund is a benefaction, a part of the salary, and that participation in it involves a kind of compulsion, albeit mild and just, he continues:
\begin{quote}
From this difference it follows, among other things, that the proposition ‘we owe present widows \textit{as much} as future ones, and conversely’, if \textit{as much} is meant to express the cash amount, is not applicable here. The fund has something of the character of shares, rising and falling under the sway of circumstances.

Increase of the widows' pensions: Mr. C. R. P. has determined the amounts of these (\textit{as also of the reduction}) by an unalterable rule.
\end{quote}
\noindent I have copied these two passages word for word because it follows from them, and in particular from the two which I have underlined twice, quite indisputably that \textsc{Sch.} understood P.'s Art. No. 3 in the sense of the first interpretation (above, p. [129]). Several others who voted subsequently, such as B., H., O., R., T., agreed with \textsc{Sch.}'s vote, although without specifying the individual questions. On the other hand, in F.'s vote to be cited [on the following p.] one recognises the second interpretation, which has also been followed in practice since 1837; and I therefore think that the wording of the rule of progression, which, as the comparison shows, passed unchanged into the rescript of 20 November 1794, can quite fairly be called \textit{unclear}.

Apart from \textsc{Sch.}, only M. mentioned our question (P.'s 3) expressly and unambiguously; but he says nothing more about it than that the \textit{establishment} of a rule for the future rise in the pension, as recommended by P., seems particularly important. One could say that, strictly speaking, this contains only an approval of the end, but not yet definitely an approval of the means proposed by P., nor a statement of the sense in which M. understood the latter. Without attaching any weight to this, however, I do not wish to leave unmentioned that \textsc{Meiners}, in his well-known work on the constitution and administration of German universities, p. 95, cites the existing rule of progression, just as \textsc{Brandes} does (see above, p. [129]), with the words: \textit{as long as} the number of widows does not exceed 15.

From P. himself we have nothing further on our question than the article given above on p. [136], Art. 3. If I may dwell a moment longer on the question of the sense in which P. himself understood it, I draw attention to the fact that the words ‘It is clear that one can do this’ admit of only two interpretations: namely, either P.'s view was that such an increase
\begin{quote}
should remain valid only as long as the number 15 had not yet been exceeded, but after it was exceeded should either cease again or be modified in a suitable way
\end{quote}
\noindent or P. considered the exceeding of the number 15 impossible, or at least so very improbable that taking such a case into account was entirely unnecessary.
